% Part: sets-functions-relations
% Chapter: arithmetization
% Section: rationals

\documentclass[../../../include/open-logic-section]{subfiles}

\begin{document}

\olfileid{sfr}{arith}{rat}
\olsection{$\Int$ توں $\Rat$ ول}

اسی ہُنّے ویکھیا کہ سادی سیٹ تھیوری ورت کے قدرتی عدداں توں صحیح عدد
کِویں بنائے جا سکدے نیں۔ ہُن اسی ویکھاں گے کہ بالکل رلدے ملدے طریقے نال
صحیح عدداں توں ناطق عدد کِویں بنائے جا سکدے نیں۔ ساڈیاں مُڈھلیاں گلّاں
ایہ نیں:
\begin{enumerate}
	\item ہر ناطق عدد نوں $\nicefrac{i}{j}$ دی صورت وچ لکھیا جا سکدا اے،
	جتھے $i$ تے $j$ دونیں صحیح عدد نیں، پر $j$ صفر نہیں۔
	\item اظہار $\nicefrac{i}{j}$ وچ کوڈ کیتی معلومات نوں ترتیب وار جوڑے
	$\tuple{ i, j}$ وچ وی اوہو جیہے طور تے کوڈ کیتا جا سکدا اے۔
\end{enumerate}
صاف طریقہ ایہ ہووے گا کہ ناطق عدداں نوں $\Int \times (\Int \setminus
\{0_\Int\})$ توں چُنے ترتیب وار جوڑے \emph{سمجھیا} جاوے۔ پر پہلے وانگ
ایہ وی ذرا حدوں ودھ سادہ ہووے گا، کیوں جو اسی چاہنے آں کہ
$\nicefrac{3}{2} = \nicefrac{6}{4}$، پر $\tuple{ 3, 2}\neq \tuple{ 6,
4}$۔ ہور عام طور تے ساڈی لوڑ ایہ اے:
\[
	\nicefrac{a}{b} = \nicefrac{c}{d} \text{ اُدوں تے صرف اُدوں جدوں } a \times d = b \times c
\]
ایہ حاصل کرن لئی اسی $\Int \times (\Int \setminus \{0_\Int\})$ اُتے
اک {ہم ارزیت دا تعلق} ایویں تعریف کردے آں:
\[
	\tuple{ a, b }\Ratequiv \tuple{ c, d} \text{ اُدوں تے صرف اُدوں جدوں }a \times d = b \times c
\]
ساڈے نوں جانچنا پینا اے کہ ایہ ہم ارزیت دا تعلق اے۔ ایہ
$\Intequiv$ والے معاملے نال بہت رلدا اے، تے اسی ایہنوں مشق لئی
چھڈدے آں۔
\begin{prob}
وکھاؤ کہ $\Ratequiv$ ہم ارزیت دا تعلق اے۔
\end{prob}
پر ایس نال اسی ایہ آکھ سکدے آں:
\begin{defn}
ناطق عدد، صحیح عدداں دیاں جوڑیاں دے $\Ratequiv$ تحت ہم ارزیت دے طبقے
نیں (جیہناں دا دوجا رکن صفر نہیں)۔ یعنی $\Rat =
\equivclass{(\Int \times (\Int\setminus \{0_\Int\}))}{\Ratequiv}$۔
\end{defn}

صحیح عدداں وانگ، اسی کجھ بنیادی عمل وی تعریف کرنا چاہنے آں۔ جتھے
$\equivrep{i,j}{\Ratequiv}$، $\Ratequiv$ تحت اوہ ہم ارزیت دا طبقہ اے
جیہدے وچ $\tuple{i, j}$، !!a{element} اے، اسی آکھدے آں:
\begin{align*}
	\equivrep{a, b}{\Ratequiv} + \equivrep{c, d}{\Ratequiv} &= \equivrep{ad + bc,  bd}{\Ratequiv}\\
	\equivrep{a, b}{\Ratequiv} \times \equivrep{c, d}{\Ratequiv} &= \equivrep{a  c, b d}{\Ratequiv}.
\intertext{ایہناں ناطق عدداں اُتے $r \leq s$ تعریف کرن لئی اسی ایہ
گل ورتدے آں کہ $r \le s$ اُدوں تے صرف اُدوں جدوں $s - r$ منفی نہیں،
یعنی $s - r$ نوں $\nicefrac{i}{j}$ دی صورت وچ لکھیا جا سکدا اے، جتھے
$i$ غیر منفی تے $j$ مثبت اے:}
	\equivrep{a, b}{\Ratequiv} \leq \equivrep{c, d}{\Ratequiv} &\text{ اُدوں تے صرف اُدوں جدوں }
	\equivrep{c, d}{\Ratequiv} - \equivrep{a, b}{\Ratequiv} =
	\equivrep{i_\Int, j_\Int}{\Ratequiv}
\end{align*}
% PNBARI-002 ماخذ دی درستی: انگریزی وضاحت اک واری فرق نوں r-s آکھدی اے؛ اگلی تعریف تے پچھلا جملہ دونیں s-r مقرر کردے نیں۔
کِسے $i \in \Nat$ تے $0 \neq j \in \Nat$ لئی۔

فیر ساڈے نوں جانچنا پینا اے کہ ایہ تعریفاں اوہو جیہا ورتاؤ کردیاں نیں
جیہو جیہا اوہناں نوں \emph{کرنا چاہیدا اے}؛ تے ایہ کم اسی
\olref[check]{sec} ول بھیج دے آں۔ پر ایہ واقعی ایویں ای کردیاں نیں!{}
آخر وچ اسی صحیح عدداں نوں ناطق عدد \emph{سمجھ کے} ورتن دا کوئی طریقہ
چاہنے آں؛ سو ہر $i \in \Int$ لئی اسی مقرر کردے آں کہ $i_\Rat =
\equivrep{i, 1_\Int}{\Ratequiv}$۔ فیر اسی \olref[check]{sec} وچ جانچدے
آں کہ ایہ سارا کم درست ورتاؤ کردا اے۔

\begin{prob}
وکھاؤ کہ $(i + j)_\Rat = i_\Rat+ j_\Rat$ تے $(i \times j)_\Rat =
i_\Rat \times j_\Rat$ تے $i \leq j \liff i_\Rat \leq j_\Rat$، ہر
$i, j \in \Int$ لئی۔
\end{prob}

\end{document}
